Un canó electrònic que dispara electrons els accelera, mitjançant un camp elèctric uniforme generat per dues plaques metàl·liques (A i B), des del repòs fins a una velocitat de $2,00\times 10^{6}\ \mathrm{m\,s^{-1}}$ (figura 1). Dins del canó, els electrons inicien el recorregut a la placa A i viatgen cap a la placa B, per on surten horitzontalment cap a la dreta per un petit orifici. Les dues plaques són paral·leles i estan separades per 4,00~cm.

{\centering \textsc{Figura 1}\par}
\figura[0.32]{fig1.png}

\begin{apartat}{1}
Calculeu la diferència de potencial entre les dues plaques i indiqueu quina placa té el potencial més alt i quina té el potencial més baix. Dibuixeu la figura 1 i representeu-hi les línies de camp elèctric entre les dues plaques.
\end{apartat}

{\centering \textsc{Figura 2}\par}
\figura[0.4]{fig2.png}

\begin{apartat}{1}
Més endavant, els electrons passen entre dues altres plaques, que generen un camp elèctric uniforme de $500\ \mathrm{N\,C^{-1}}$ vertical cap amunt (figura 2). Calculeu l'acceleració dels electrons quan estiguin sota l'acció d'aquest camp elèctric i les dues components de la velocitat en sortir del recinte on hi ha el camp elèctric.
\end{apartat}

\dades{%
  $|e| = 1,6\times 10^{-19}\ \mathrm{C}$.\\
  $m_\mathrm{e} = 9,11\times 10^{-31}\ \mathrm{kg}$.}
\nota{Considereu negligible el camp gravitatori.}
